% Fichier généré automatiquement par tools/generate_corrections.py.
% Source : GEO/GEO-01/14_Sympact/14_Sympact.tex
% Statut : partiel (4/5 question(s) renseignée(s), 0 reprise(s)).
\begin{corrigebox}[Corrigé partiel extrait de la banque]
\CorrigeQuestion{1}
\begin{center}
\includegraphics[width=.4\linewidth]{14_02_c}
\end{center}
\CorrigeQuestion{2}
\begin{center}
\includegraphics[width=.5\linewidth]{14_01_c}
\end{center}
\CorrigeQuestion{3}
\CorrigeACompleter{Aucun élément de réponse n'est présent dans la source d'origine pour cette question.}
\CorrigeQuestion{4}
Dans cas, dans le pire des cas, $A$, $B$ et $C$ sont alignés (avec $B$ au-dessus de $C$). Il faut donc $AB = AC+CB = \SI{160}{mm}$.
\CorrigeQuestion{5}
Comme je suis paresseux, j'ai réalisé la construction avec geogebra. On mesure \SI{160,8}{\degres}.
\begin{marginfigure}
\rotatebox{90}{\includegraphics[width=.35\linewidth]{14_03_c}}
\end{marginfigure}

\begin{enumerate}
\item .
\item .
\item .
\item \SI{160}{mm}.
\item \SI{160,8}{\degres}.
\end{enumerate}
\end{corrigebox}
